% My first latex document

\documentclass{article}
\usepackage[utf8]{inputenc}
\usepackage{graphicx}
\usepackage[ngerman]{babel}

\title{Erstes LaTeX Dokument}
\date{\today}
\author{Simon Struck}
\begin{document}
	\maketitle
	\newpage
	\tableofcontents
	\newpage
	\section{One}
	\subsection{One.One}
	\paragraph{Paragraph}
	Paragraph text
	Paragraph text
	Paragraph text
	Paragraph text
	\\
	\\
	Paragraph text
	Paragraph text
	\\
	\subparagraph[Short sub paragraph]{Long subparagraph} Long subparagraph Long subparagraph Long subparagraph Long subparagraph Long subparagraph Long subparagraph Long subparagraph Long subparagraph 
	\newpage
	\section{Add images}
	\begin{figure}
		\centering
		\includegraphics[width=0.5\linewidth,height=3cm]{img/logo.png}
		\caption{Logo der Hochschule}
		\label{fig:logo}
	\end{figure}
	Auf Abbildung \ref{fig:logo} sieht man das Logo der Hs Emden/Leer.
	Niemals sollte \textbackslash label vor \textbackslash caption erstellt werden.
	
	
	\subsection{Add image to custom places}
	
	\begin{figure}[b]
		\centering
		\includegraphics[width=0.5\linewidth]{img/test.jpg}
		\caption{Testbild}
		\label{fig:test}
		\end{figure}
	Die Abbildung \ref{fig:test} wurde mit dem Befehl [b] nach unten auf die Seite verschoben.
	Mit [h] oder [h!] kann man das Bild auch zwingen im Text zu platzieren.
	
	Das Paket "placeins" kann Bilder noch besser in Fließtext einbauen.
	\newpage
	\appendix
	\listoffigures
\end{document}